\section{Periodic sequences and fundamental periods (2.37--2.42)}

\Prob{2.37\quad One period of the $N$-periodic extension of the Problem-2.3 sequences}

The $N$-periodic extension is
$\tilde{s}[n]=\sum_{k}s[n+kN]$, so one period is obtained by folding every sample of $s$ into
the window $0\le n\le N-1$ at position $n\bmod N$ and \emph{adding} whatever lands on the
same slot.

\textbf{(a) $N=6$.} With $x[n]$ at $n=-1,\dots,5$, the sample $x[-1]=3$ folds to slot
$-1\bmod 6=5$, where $x[5]=-3$ already sits: $3+(-3)=0$.
\[
\begin{array}{r|rrrrrr}
n & 0&1&2&3&4&5\\\hline
\tilde{x}[n] & -4&-2&5&0&1&0\\
\tilde{y}[n] & 3&6&-2&3&-1&0\\
\tilde{w}[n] & 0&1&-5&5&3&-4
\end{array}
\]
(For $y$: $y[-3]=-5\to$ slot 3 joins $y[3]=8$ giving $3$; $y[-2]=-1\to$ slot 4 gives $-1$;
$y[-1]=0\to$ slot 5 gives $0$. For $w$: indices $3\dots9$ fold to slots $3,4,5,0,1,2,3$, so
slot 3 receives $w[3]+w[9]=-1+6=5$.)

\textbf{(b) $N=8$.}
\[
\begin{array}{r|rrrrrrrr}
n & 0&1&2&3&4&5&6&7\\\hline
\tilde{x}[n] & -4&-2&5&0&1&-3&0&3\\
\tilde{y}[n] & 3&6&-2&8&0&-5&-1&0\\
\tilde{w}[n] & -5&6&0&-1&3&-4&0&1
\end{array}
\]
Here $x$ and $y$ (length 7) suffer no overlap --- only $w$'s indices $3\dots 9$ wrap, with
$w[8]=-5$ landing in slot 0 and $w[9]=6$ in slot 1.

Sanity check: for every $N$, $\sum_{n=0}^{N-1}\tilde{s}[n]=\sum_n s[n]$. Indeed
$\sum\tilde x=0=\sum x$, $\sum\tilde y=9=\sum y$, $\sum \tilde w=0=\sum w$ in both tables.

\Prob{2.38\quad Identify $A,\omega_o,\phi$ from one period of $\tilde{x}[n]=A\sin(\omega_o n+\phi)$}

If one period contains $N$ samples then $\omega_o = 2\pi r/N$ for some integer $r$ coprime
with $N$; the remaining two unknowns follow from two samples.

\textbf{(a)} $\{1,\,1,\,-1,\,-1\}$: $N=4\Rightarrow\omega_o=\pi/2$ (taking $r=1$).
Then $A\sin\phi=x[0]=1$ and $A\sin(\pi/2+\phi)=A\cos\phi=x[1]=1$. Dividing,
$\tan\phi=1\Rightarrow\phi=\pi/4$, and $A=1/\sin(\pi/4)=\sqrt2$.
\[
\boxed{\;A=\sqrt{2},\quad \omega_o=\frac{\pi}{2},\quad \phi=\frac{\pi}{4}.\;}
\]
Check: $\sqrt2\sin(3\pi/4)=1$ \checkmark, $\sqrt2\sin(\pi+\pi/4)=-1$ \checkmark.

\textbf{(b)} $\{0.5,\,-0.5,\,0.5,\,-0.5\}$: the sequence actually repeats every $2$ samples,
so $N=2$ and $\omega_o=\pi$. Since
$\sin(\pi n+\phi)=(-1)^{n}\sin\phi$, we need $A\sin\phi=0.5$; the natural normalized choice
is $\phi=\pi/2$:
\[
\boxed{\;A=0.5,\quad \omega_o=\pi,\quad \phi=\frac{\pi}{2}.\;}
\]
(Only the product $A\sin\phi$ is determined by the data --- any $(A,\phi)$ with
$A\sin\phi=0.5$ reproduces the sequence. This degeneracy is intrinsic to $\omega_o=\pi$.)

\textbf{(c)} $\{0,\,0.5878,\,-0.9511,\,0.9511,\,-0.5878\}$: $N=5$. Since $x[0]=0$, take
$\phi=0$. Then $x[1]=A\sin\omega_o=0.5878$. Among $\omega_o=2\pi r/5$,
$\sin(4\pi/5)=0.5878$, so $r=2$ and $A=1$:
\[
\boxed{\;A=1,\quad \omega_o=\frac{4\pi}{5}=0.8\pi,\quad \phi=0.\;}
\]
Check: $\sin(8\pi/5)=-0.9511$ \checkmark, $\sin(12\pi/5)=\sin(2\pi/5)=0.9511$ \checkmark,
$\sin(16\pi/5)=\sin(6\pi/5)=-0.5878$ \checkmark.

\textbf{(d)} $\{2,\,0,\,-2,\,0\}$: $N=4\Rightarrow \omega_o=\pi/2$.
$x[1]=A\cos\phi=0\Rightarrow \phi=\pm\pi/2$; with $\phi=\pi/2$, $x[0]=A\sin(\pi/2)=A=2$:
\[
\boxed{\;A=2,\quad \omega_o=\frac{\pi}{2},\quad \phi=\frac{\pi}{2}.\;}
\]
(Equivalently $x[n]=2\cos(\pi n/2)$.)

\Prob{2.39\quad Fundamental periods}

A sequence $\cos(\omega_o n+\phi)$ (or $e^{j\omega_o n}$) is periodic iff
$\omega_o N=2\pi r$ for integers $N>0,r$, i.e.\ iff $\omega_o/2\pi$ is rational; the
\emph{fundamental} period is the smallest such $N$:
\[
N=\frac{2\pi r}{\omega_o}\Bigg|_{\text{smallest integer }N} .
\]
For a sum of sinusoids, $N=\mathrm{LCM}$ of the individual fundamental periods.

\textbf{(a)} $\tilde{x}_a[n]=e^{j0.25\pi n}$: $\omega_o=0.25\pi$, $N=\dfrac{2\pi r}{0.25\pi}=8r$.
\[
\boxed{\;N=8 \quad (r=1).\;}
\]

\textbf{(b)} $\tilde{x}_b[n]=\cos(0.6\pi n+0.3\pi)$: $N=\dfrac{2\pi r}{0.6\pi}=\dfrac{10r}{3}$,
smallest integer at $r=3$.
\[
\boxed{\;N=10 .\;}
\]
(The phase never affects the period.)

\textbf{(c)} $\tilde{x}_c[n]=\operatorname{Re}\bigl(e^{j\pi n/8}\bigr)+\operatorname{Im}\bigl(e^{j\pi n/5}\bigr)
=\cos(\pi n/8)+\sin(\pi n/5)$:
\[
\omega_1=\tfrac{\pi}{8}\Rightarrow N_1=\frac{2\pi}{\pi/8}=16,\qquad
\omega_2=\tfrac{\pi}{5}\Rightarrow N_2=\frac{2\pi}{\pi/5}=10 .
\]
\[
\boxed{\;N=\mathrm{LCM}(16,10)=\frac{16\cdot 10}{\gcd(16,10)}=\frac{160}{2}=80 .\;}
\]

\textbf{(d)} $\tilde{x}_d[n]=6\sin(0.15\pi n)-\cos(0.12\pi n+0.1\pi)$:
\[
N_1=\frac{2\pi r}{0.15\pi}=\frac{40r}{3}\Rightarrow N_1=40\ (r=3),\qquad
N_2=\frac{2\pi r}{0.12\pi}=\frac{50r}{3}\Rightarrow N_2=50\ (r=3).
\]
\[
\boxed{\;N=\mathrm{LCM}(40,50)=\frac{2000}{10}=200 .\;}
\]

\textbf{(e)} $\tilde{x}_e[n]=\sin(0.1\pi n+0.75\pi)-3\cos(0.8\pi n+0.2\pi)+\cos(1.3\pi n)$:
\[
N_1=\frac{2\pi r}{0.1\pi}=20r\Rightarrow 20,\qquad
N_2=\frac{2\pi r}{0.8\pi}=2.5r\Rightarrow 5\ (r=2),\qquad
N_3=\frac{2\pi r}{1.3\pi}=\frac{20r}{13}\Rightarrow 20\ (r=13).
\]
\[
\boxed{\;N=\mathrm{LCM}(20,5,20)=20 .\;}
\]

\Prob{2.40\quad Fundamental period of $\tilde{x}[n]=A\sin(\omega_o n)$}

Solve $N=2\pi r/\omega_o$ for the smallest positive integer $N$; equivalently, write
$\omega_o/2\pi = r/N$ in lowest terms.

\begin{center}
\renewcommand{\arraystretch}{1.5}
\begin{tabular}{@{}cllc@{}}
\toprule
& $\omega_o$ & $\dfrac{\omega_o}{2\pi}$ in lowest terms & $N$\\
\midrule
(a) & $0.3\pi$   & $\dfrac{0.3}{2}=\dfrac{3}{20}$    & $\boxed{20}$\\
(b) & $0.48\pi$  & $\dfrac{0.48}{2}=\dfrac{6}{25}$   & $\boxed{25}$\\
(c) & $0.45\pi$  & $\dfrac{0.45}{2}=\dfrac{9}{40}$   & $\boxed{40}$\\
(d) & $0.525\pi$ & $\dfrac{0.525}{2}=\dfrac{21}{80}$ & $\boxed{80}$\\
(e) & $0.75\pi$  & $\dfrac{0.75}{2}=\dfrac{3}{8}$    & $\boxed{8}$\\
(f) & $0.75\pi$  & $\dfrac{0.75}{2}=\dfrac{3}{8}$    & $\boxed{8}$\\
\bottomrule
\end{tabular}
\end{center}
\Mis parts (e) and (f) are printed with the same value $0.75\pi$; they therefore have the
same answer. (In the electronic 4th-edition file this problem uses the list
$0.6\pi,\,0.28\pi,\,0.45\pi,\,0.55\pi,\,0.65\pi$, giving $N=10,\,50,\,40,\,40,\,40$.)

Notice how sensitive $N$ is to $\omega_o$: $0.45\pi$ and $0.525\pi$ differ by less than
$17\%$ yet their periods are $40$ and $80$.

\Prob{2.41\quad Period of $\tilde{x}_1[n]=\sin(0.06\pi n)$ and other sequences with the same period}

\[
\frac{\omega_o}{2\pi}=\frac{0.06\pi}{2\pi}=\frac{3}{100}\quad\text{(already in lowest terms)}
\qquad\Longrightarrow\qquad \boxed{\;N=100 .\;}
\]

Any sinusoid $\sin(\omega n+\phi)$ has fundamental period exactly $100$ provided
$\omega=\dfrac{2\pi r}{100}=\dfrac{\pi r}{50}$ with $\gcd(r,100)=1$. Choosing $r=7$ and
$r=13$:
\[
\boxed{\;\tilde{x}_2[n]=\sin(0.14\pi n),\qquad \tilde{x}_3[n]=\sin(0.26\pi n).\;}
\]
Both are distinct from $\tilde{x}_1[n]$ and from each other, and both have fundamental period
$100$. (Other valid choices: $r=9,11,17,\dots$, or a cosine, or any phase shift. Note that
$r=5$ would \emph{not} work: $\gcd(5,100)=5$ gives $N=20$.)

\Prob{2.42\quad Fundamental period and average power of periodic sequences}

For a single sinusoid $A\cos(\omega_o n+\phi)$ with fundamental period $N>2$, the average
power is $A^{2}/2$ (Problem 2.36(e)). For a \emph{sum} of sinusoids with distinct
frequencies in $(0,\pi)$, the cross terms average to zero over the common period, so the
powers add:
\[
\mathcal{P}=\sum_i \frac{A_i^{2}}{2}.
\]

\textbf{(a)} $\tilde{x}_1[n]=5\cos(\pi n/3)$: $\dfrac{\omega_o}{2\pi}=\dfrac{1}{6}$.
\[
\boxed{\;N=6,\qquad \mathcal{P}=\tfrac{5^{2}}{2}=12.5 .\;}
\]

\textbf{(b)} $\tilde{x}_2[n]=2\cos(2\pi n/5)$: $\dfrac{\omega_o}{2\pi}=\dfrac15$.
\[
\boxed{\;N=5,\qquad \mathcal{P}=\tfrac{2^{2}}{2}=2 .\;}
\]

\textbf{(c)} $\tilde{x}_3[n]=2\cos(2\pi n/7)$: $\dfrac{\omega_o}{2\pi}=\dfrac17$.
\[
\boxed{\;N=7,\qquad \mathcal{P}=2 .\;}
\]

\textbf{(d)} $\tilde{x}_4[n]=3\cos(5\pi n/7)$: $\dfrac{\omega_o}{2\pi}=\dfrac{5}{14}$,
$\gcd(5,14)=1$.
\[
\boxed{\;N=14,\qquad \mathcal{P}=\tfrac{3^{2}}{2}=4.5 .\;}
\]

\textbf{(e)} $\tilde{x}_5[n]=4\cos(2\pi n/5)+3\cos(3\pi n/5)$:
$N_1=5$ and $\dfrac{3\pi/5}{2\pi}=\dfrac{3}{10}\Rightarrow N_2=10$, so
$N=\mathrm{LCM}(5,10)=10$. The two frequencies $0.4\pi$ and $0.6\pi$ are distinct and
neither is $0$ nor $\pi$, and $0.4\pi\ne 2\pi-0.6\pi$, so the cross term averages to zero
over $N=10$:
\[
\frac{1}{10}\sum_{n=0}^{9}\cos(0.4\pi n)\cos(0.6\pi n)
=\frac{1}{20}\sum_{n=0}^{9}\bigl[\cos(\pi n)+\cos(0.2\pi n)\bigr]=0 .
\]
\[
\boxed{\;N=10,\qquad \mathcal{P}=\tfrac{16}{2}+\tfrac{9}{2}=12.5 .\;}
\]

\textbf{(f)} $\tilde{x}_6[n]=4\cos(5\pi n/3)+3\cos(3\pi n/5)$:
$\dfrac{5\pi/3}{2\pi}=\dfrac{5}{6}\Rightarrow N_1=6$; $N_2=10$ as above; so
$N=\mathrm{LCM}(6,10)=30$. Again the frequencies are distinct modulo the folding
$\omega\leftrightarrow 2\pi-\omega$ (note $5\pi/3$ aliases to $\pi/3$, which is still
$\ne 3\pi/5$), so the powers add:
\[
\boxed{\;N=30,\qquad \mathcal{P}=\tfrac{16}{2}+\tfrac{9}{2}=12.5 .\;}
\]
